\renewcommand{\thesubsection}{(\alph{subsection})}

\section{Problem 1: Cart-Pole with Actuator Dynamics: DDP and iLQR}
\subsection*{Prerequisite: cart-Pole Dynamics}
The dynamic equation of cart-pole system is as follows:
\begin{align}
\ddot{x} &= \frac{\ell m_p \sin(\theta)\,\dot{\theta}^{\,2} + F + m_p g \cos(\theta)\sin(\theta)}
{
    m_c + m_p\left(1-\cos^2(\theta)\right)
}\\
\ddot{\theta} &= -\frac{
    \ell m_p \cos(\theta)\sin(\theta)\,\dot{\theta}^{\,2}
    + F\cos(\theta)
    + (m_c + m_p)g\sin(\theta)
}{
    \ell m_c + \ell m_p\left(1-\cos^2(\theta)\right)
}
\end{align}

\subsection{ }
Denote $\mathbf{x} = [x, \theta, \dot{x}, \dot{\theta}, F]$ as the augumented state vector. The dynamics of cart-pole system can be written as follows:

\begin{equation}
    \dot{\mathbf{x}} = 
\begin{bmatrix}
    \dot{x}\\
    \dot{\theta}\\
    \ddot{x}\\
    \ddot{\theta}\\
    \dot{F}
\end{bmatrix}
= 
\begin{bmatrix}
    \dot{x}\\
    \dot{\theta}\\
    \frac{\ell m_p \sin(\theta)\,\dot{\theta}^{\,2} + F + m_p g \cos(\theta)\sin(\theta)}
{
    m_c + m_p\left(1-\cos^2(\theta)\right)
}\\
    -\frac{
    \ell m_p \cos(\theta)\sin(\theta)\,\dot{\theta}^{\,2}
    + F\cos(\theta)
    + (m_c + m_p)g\sin(\theta)
}{
    \ell m_c + \ell m_p\left(1-\cos^2(\theta)\right)
}\\
    \frac{u-F}{\tau_a}
\end{bmatrix}
= \boldsymbol{f}(\mathbf{x},u)
\end{equation}

\begin{itemize}
    \item Forward Euler Discretization
\end{itemize}

Given the continuous-time nonlinear dynamics

\begin{equation}
\dot{\mathbf{x}} = \boldsymbol{f}(\mathbf{x},u),
\end{equation}

the Forward Euler method with sampling time \(h\) approximates the state evolution as

\begin{equation}
\mathbf{x}_{k+1}
=
\mathbf{x}_k
+
h \boldsymbol{f}(\mathbf{x}_k,u_k).
\end{equation}

For the augmented cart-pole state the discrete-time dynamics become
\begin{align}
x_{k+1}
&=
x_k + h \dot{x}_k, \\
\theta_{k+1}
&=
\theta_k + h \dot{\theta}_k, \\
\dot{x}_{k+1}
&=
\dot{x}_k + h \ddot{x}_k, \\
\dot{\theta}_{k+1}
&=
\dot{\theta}_k + h \ddot{\theta}_k, \\
F_{k+1}
&=
F_k + h \frac{u_k - F_k}{\tau_a}.
\end{align}

Therefore, the discrete nonlinear dynamics can be written compactly as

\begin{equation}
\mathbf{x}_{k+1}
=
F_d(\mathbf{x}_k,u_k)
=
\mathbf{x}_k + h \boldsymbol{f}(\mathbf{x}_k,u_k).
\end{equation}

As an alternative, a higher-order numerical integration method such as the fourth-order Runge--Kutta (RK4) scheme may be used to improve integration accuracy.

\subsection{}
Define $\mathbf{X} = [\mathbf{x}_1, ... \mathbf{x}_N]$ and $ \mathbf{U} = [u_0,...,u_{N-1}]$ as matrices of state and control trajectory. Then a cost function for the swing-up task is as follows:
\begin{equation}
    J(\mathbf{X},\mathbf{U}) = \frac{1}{2} (\mathbf{x}_N -\mathbf{x}_{ref})^T\mathbf{Q}_N (\mathbf{x}_N -\mathbf{x}_{ref}) + \frac{1}{2}\sum_{k = 0}^{N-1}\left((\mathbf{x}_k -\mathbf{x}_{ref})^T \mathbf{Q} (\mathbf{x}_k -\mathbf{x}_{ref}) + Ru_k^2 \right)
\end{equation}
where
\begin{equation}
    l_N({\mathbf{x}_N}) = \frac{1}{2} (\mathbf{x}_N -\mathbf{x}_{ref})^T\mathbf{Q}_N (\mathbf{x}_N -\mathbf{x}_{ref})
\end{equation}
is the terminal cost and 
\begin{equation}
    \sum_{k = 0}^{N-1}l_k (\mathbf{x}_k, u_k)  = \frac{1}{2}\sum_{k = 0}^{N-1}\left((\mathbf{x}_k -\mathbf{x}_{ref})^T \mathbf{Q} (\mathbf{x}_k -\mathbf{x}_{ref}) + Ru_k^2 \right)
\end{equation}
is the running cost. Define the value function:
\begin{equation}
    V(\mathbf{x}_k) = \min_{u_k,u_{k+1}, \ldots , u_{N-1}}\left(\sum_{t = k}^{N-1}l_k (\mathbf{x}_t, u_t) +   l_N(\mathbf{x}_N) \right)
\end{equation}
The Bellman optimality condition provides a recursion on $V(\mathbf{x}_k)$ as follows:
\begin{align}
    V(\mathbf{x}_k)  &= \min_{u_k}\left(l_k(\mathbf{x}_k,u_k) + V(\mathbf{x}_{k+1})\right)\\
    &= \min_{u_k}\left(l_k(\mathbf{x}_k,u_k) + V_{k+1}(F_d(\mathbf{x}, u))\right)
\end{align}
\subsection{}
Define the local action-value function $Q(\delta\mathbf{x}_k, \delta u_k)$ as follows:
\begin{equation}
    Q(\delta\mathbf{x}_k, \delta u_k) = l_k(\mathbf{x}_k + \delta\mathbf{x}_k,u_k+\delta u_k) + V_{k+1}(F_d(\mathbf{x}_k + \delta\mathbf{x}_k,u_k+\delta u_k))
    \label{eq: action value}
\end{equation}
Then we can approximate $Q(\delta\mathbf{x}_k, \delta u_k)$ in a quadratic form by expanding to 2-nd order of its Taylor expansion:
\begin{align}
    \label{eq: Q value quadratic}
    Q(\delta\mathbf{x}_k, \delta u_k) \approx Q(\mathbf{0},0) + Q_{x,k}^T \delta \mathbf{x} + Q_{u,k}^T \delta u_k &+ \frac{1}{2}\left(\delta \mathbf{x}_k^T Q_{xx,k}\delta\mathbf{x}_k + \delta u_k^T Q_{uu,k}\delta u_k\right)\\ 
   &+ \delta u_k^T Q_{ux,k}\delta \mathbf{x}_k \nonumber   
\end{align}
where from definition of \ref{eq: action value} we can determine:
\begin{align*}
Q_k(0,0) &= l(\mathbf{x}_k, u_k) + V_{k+1}(\mathbf{x}_{k+1}), \\
Q_{x} &= l_{x} + F_{dx}^{T} V_{x,k+1}, \\
Q_{u} &= l_{u} + F_{d,u}^{T} V_{x,k+1}, \\
Q_{xx} &= l_{xx} + F_{dx}^{T} V_{xx,k+1} F_{dx} + V_{x,k+1}^{T} F_{dxx}, \\
Q_{uu} &= l_{uu} + F_{d,u}^{T} V_{xx,k+1} F_{du} + V_{x,k+1}^{T} F_{duu}, \\
Q_{ux} &= l_{ux} + F_{du}^{T} V_{xx,k+1} F_{dx} + V_{x,k+1}^{T} F_{dux}.
\end{align*}
where the subscription $k+1$ indicate the value taken at time $k+1$. Otherwise we evaluate at time $k$.
\subsection{}
Now we minimize $Q(\delta \mathbf{x}_k, \delta u_k)$ in \ref{eq: Q value quadratic} over $\delta u_k$, we have
\begin{align}
    \delta u_k &= \arg \min_{\delta u_k}{Q(\delta \mathbf{x}_k, \delta u_k)} = \arg \min_{\delta u_k} \frac{1}{2}\delta u_k^T Q_{uu,k}\delta u_k  + Q_{u,k}^T \delta u_k + \delta u_k^T Q_{ux,k}\delta \mathbf{x}_k\\
    &= -Q_{uu}^{-1}
\end{align}


